\documentclass[aspectratio=169]{beamer}
\usetheme{Madrid}
\title{Fourier Series}
\subtitle{Building signals from simple waves}
\author{Beamer Deck Maker}
\date{}

\begin{document}

\begin{frame}
  \titlepage
\end{frame}

\begin{frame}{The idea}
  A Fourier series writes a periodic function as a sum of harmonics.
  \[
    f(x) \sim \frac{a_0}{2}
    + \sum_{n=1}^{\infty} \bigl(a_n\cos(nx) + b_n\sin(nx)\bigr)
  \]
  \begin{itemize}
    \item Each harmonic oscillates at an integer multiple of the base frequency.
    \item The coefficients tell us how much each harmonic contributes.
  \end{itemize}
\end{frame}

\begin{frame}{A square wave}
  The unit square wave that is positive on $(0,\pi)$ and negative on
  $(-\pi,0)$ contains only odd sine harmonics.
  \[
    f(x) \sim \frac{4}{\pi}\left(\sin x + \frac{\sin 3x}{3}
    + \frac{\sin 5x}{5} + \cdots\right)
  \]
  Adding harmonics makes the flat sections more accurate and the transitions sharper.
\end{frame}

\end{document}
